\documentclass{article}
\begin{document}
\section*{Binary native route and flat-chart comparison}
With the four local binary and four nullary middle-four distributor
diagrams from the first tranche, and ordinary distributor naturality,
\texttt{native\_binary\_on\_chart} compares independently composed
distributors on the addressed row-major branch chart. The four new
reductions recover the literal factorization in each of the two inner
and two outer stages. The continuation maps and parent merges remain
whole-sum arrows. The reverse-function model instantiates all these
premises on actual finite native homs.

The joint native carrier preimage is proved equivalent to the two
dependent preimages, each including its raw parent and child address.
A typed adjacent-swap program and a nested branch tree now construct an
invertible structural comparison from a flat sum, preserving and
reflecting the equations after change of chart. These chart comparisons
are natural in every whole coefficient leaf arrow. In the reverse-function
model all swap programs with the same endpoints agree when the source
list has no duplicate raw addresses; this uses recovery of the tagged
token, not equality of its type or value. An empty list contributes the
additive unit; an empty block retains its chart position.

The general structural identification of the original chosen flat native
maps with the two row-major factorizations is not yet proved. To finish
it, compare the actual ordered branch histories in both independent-first
and dependent-first expansions and derive the equality of the resulting
mixed sum/tensor structural maps, including absorber cases. The generic
chart theorem transports along a supplied addressed permutation but does
not prove that mixed comparison. The abstract derivation of the eight
middle-four distributor diagrams from a standard rig presentation,
global dependent functoriality, and all-arity coherence remain separate.
\end{document}
